% Français 9117, batch 5 — Gallica views 81–100.
% Pass: Opus 5 (claude-opus-5), 2026-09-14 — first pass, unchecked against
% the views by a human.
%
% What the twenty views hold. Twenty consecutive written pages of one prose
% manuscript in Germain's hand — the « Considérations générales sur l'état
% des sciences et des lettres » — ten leaves, recto and verso, no blank, no
% mathematics. It is a working copy, not a fair copy: every page carries
% deletions and interlinear rewrites, several paragraphs open with a square
% bracket of the author's, one (v. 85) is marked « Variante » in the margin,
% and one (v. 93) is cancelled by a large cross. The hand is her regular
% cursive, faster and blotted on views 88, 90 and 96.
%
% View 81 opens mid-sentence, continuing the last line of view 80
% (« … dans des recherches où l'objet de / les études étoient placés hors
% d'elle »). View 100 ends on a complete sentence with space left below; the
% argument goes on to the next batch.
%
%   v. 81      The imitation of ancient literature, its fictions faded in a
%              nation that did not invent them; a new school seeking a
%              literature of our own. The invasion of mathematical forms
%              into works whose nature cannot reach their exactness.
%   v. 82–83   The style of the higher mathematical works, and its grotesque
%              copies in common language; figures used where they express no
%              real value; nations, like beginners in serious study, trusting
%              blindly the doctrines they are taught.
%   v. 84–85   Pedantry, once the scholar's fault, now that of whole masses;
%              the young scorning the old playful tone. « Ne nous alarmons
%              pas » : politics will soon adopt the forms proper to it; then,
%              marked « Variante », the claim that the vaunted progress will
%              prove the development of ideas already in the predecessors.
%   v. 86–88   Letters have lost their lustre, but a purer day will return:
%              all branches of knowledge must develop in parallel, as in the
%              education of a pupil. Analogy will reveal the type of the true
%              in every subject; the moral and political sciences will join
%              the exact ideas; and the arts, which differ in their means,
%              all draw their power from imitation, which derives from the
%              sense of analogy.
%   v. 89–90   Laws of being relating a « module » to its subject; eloquence,
%              music and painting proceed alike — calm openings, gradation,
%              the crepuscule before sunrise, the modern ending by diminished
%              intensity of sound.
%   v. 91–93   One condition for perceiving these identities: familiarity
%              with each « module ». Why musical taste is left untrained
%              while literary taste is cultivated; the ear can be formed.
%              Music as a true language — « elle a ses phrases, ses périodes,
%              ses repos, ses hardiesses » — and the ancients' account of the
%              modes.
%   v. 94–100  Music is « toute métaphysique »: it expresses only feelings.
%              The composer, like the orator, has a dominant idea; grace and
%              gaiety, which differ by the manner of saying in speech and by
%              tempo in music; melancholy and despair, one cause, two
%              aspects; in short pieces tempo suffices, in large ones the
%              characters must be distinct — monotony and soft syllables or
%              notes of equal value for dejection, hard notes and energetic
%              expressions for violent affection, and the unexpected word or
%              note that makes the effect.
%
% Seams inside the batch. View 95 opens with the end of the sentence broken
% off at the foot of view 94 (« mais un charme [toujours] / soutenu »). View
% 97 opens « les différens genres, liés entre eux … » after a complete
% sentence at the foot of view 96: the join does not construe, and it is
% noted at the spot. Elsewhere the text runs on from page to page.
%
% The numbering on the leaves, as read. Two series.
%  (1) In ink, in the top left corner of every page, struck through: 71 at
%      view 81, then one per view to 90 at view 100 (71, 72, 73, 74, 75, 76,
%      77, 78, 79, 80, 81, 82, 83, 84, 85, 86, 87, 88, 89, 90). The digits of
%      75 (v. 85) and 84 (v. 94) are half lost under the stroke and read with
%      some doubt; the others are clear. This looks like an earlier
%      pagination of the manuscript, cancelled; it is not written into the
%      text.
%  (2) In pencil, thin and grey, in the top right corner of the rectos only:
%      40 (v. 81), 41 (v. 83), 42 (v. 85), 43 (v. 87), 44 (v. 89), 45 (v. 91),
%      46 (v. 93), 47 (v. 95), 48 (v. 97), 49 (v. 99). This is taken as the
%      library's foliation and written \folio{40r} … \folio{49r} on those
%      views. The versos carry no pencil number, and no \folio{} is written
%      on them: a verso's folio is not read from the leaf.
%  Beside the pencil 45 on view 91 stands a small struck ink figure, perhaps
%  « 4 ». Further marks not in her hand, transcribed in notes: a pencil
%  « Ch. 6 » and « avenir de » at view 84, a pencil word read « Double » at
%  the head of view 89, the word « Chévenon » (reading doubtful) in a larger
%  heavier hand at view 92, and « Georges » with a paraph at the head of view
%  99.
%
% Notation and conventions. No mathematics. Her square brackets opening a
% paragraph are kept as « [ » in the text; none is closed on these views.
% Her « module » (v. 89, 91, 92) is set \uncertain{} each time: the word is
% clear in shape but its sense here is not settled. Spelling of the page is
% kept — « étoit », « connoître », « tems », « systêmes », « poëte »,
% « hazardés », « appercevoir », « acquerreroient », « alternement »,
% « employera » — and slips are left as written (« d'homme capables », v.
% 92; « La même chose est vrai », v. 97; « peut éloigner ou même fasse »,
% v. 96). Additions written between the lines are placed where the sense
% puts them, and the notes say which they are. \uncertain{} renders as an
% underline in the PDF, as \emph{} does; she underlines nothing here.
%
% What could not be read. \ill{} stands for: a blackened word in the struck
% « contient … a portant » (v. 82); a struck word after « à l'abri » and
% another after « et qui d'ailleurs » (v. 83); the pencil words after
% « Avenir de » (v. 84); a struck word after « le vague et » (v. 85); a
% struck word after « leurs progrès » (v. 86); a struck word in « qu'il sait
% employer … que constituans » (v. 88); two struck words at view 90 (after
% « qu'elle » and after « fatiguant »); a struck word after « fait » (v. 91);
% and a struck word before « été simplement gracieuse » (v. 96). Readings
% offered with doubt are marked \uncertain{}. Nothing was guessed to fill a
% gap.
%
% Nothing here was checked against any published transcription.
\documentclass[11pt,a4paper]{article}
% The preamble shared by every transcription of the mathematicians' manuscripts in Gallica.
%
% It rests on one idea: **uncertainty compounds, it does not resolve.** A
% transcription that smooths over a doubtful word has destroyed the only thing
% it was made for — a reader must be able to tell, at every point, what was on
% the page and what was supplied. The apparatus macros below are therefore the
% critical apparatus, not decoration.
%
% The same commands are recognised by `scripts/render.mjs`, which produces the
% site's reading view, and by `scripts/tei.mjs`. A macro added here must be
% added there too, or rendering fails loudly — which is the wanted behaviour.
%
% Comments here are English, like the rest of the repository. The documents
% this preamble sets are French, and that is deliberate: see the skills.

\ProvidesPackage{germain}[2026/09/15 Transcriptions of mathematicians' manuscripts in Gallica]

\usepackage{amsmath,amssymb,amsthm}
\usepackage{mathrsfs}
% Kept from the parent preamble so the renderer's diagram support stays
% available; these manuscripts draw no commutative diagrams, and nothing here uses it.
\usepackage{tikz-cd}
\usepackage[english,french]{babel}
\usepackage{fontspec}

% --- Greek ------------------------------------------------------------------
% The text face stays fontspec's default, Latin Modern, which has no Greek:
% XeTeX drops every glyph it lacks with a line in the log and nothing on the
% page, and the Greek words the manuscripts quote vanished from the PDFs. So
% the Greek and Greek Extended blocks get a XeTeX character class of their own,
% and entering or leaving a run of it switches the family to CMU Serif — the
% same Computer Modern design, with polytonic Greek — and back. Only the
% family changes: a Greek word in \textit or \textbf keeps its shape and series.
% The transitions to and from every other class carry the tokens that class
% has towards an ordinary letter, so French spacing before « ; » or inside
% « » still holds after a Greek word. No grouping, so no brace or \endgroup
% can come out unbalanced; maths is left alone, since XeTeX inserts nothing
% there. Kept identical in `exercices.sty`.
\makeatletter
\newfontfamily\germainGreekfont{cmun}[
  Extension=.otf, NFSSFamily=germaingreek,
  UprightFont=*rm, ItalicFont=*ti, SlantedFont=*sl,
  BoldFont=*bx, BoldItalicFont=*bi, BoldSlantedFont=*bl]
\newXeTeXintercharclass\germain@greek
\def\germain@greekrange#1#2{%
  \@tempcnta=#1\relax
  \loop
    \XeTeXcharclass\@tempcnta=\germain@greek
    \ifnum\@tempcnta<#2\advance\@tempcnta\@ne\repeat}
\germain@greekrange{"0370}{"03FF}
\germain@greekrange{"1F00}{"1FFF}
\def\germain@greekprev{\rmdefault}
\def\germain@greekin{%
  \edef\germain@greekprev{\f@family}\fontfamily{germaingreek}\selectfont}
\def\germain@greekout{\fontfamily{\germain@greekprev}\selectfont}
\def\germain@greekjoin{%
  \XeTeXinterchartokenstate=\@ne
  \@tempcnta=\z@
  \loop
    \ifnum\@tempcnta=\germain@greek\else
      \XeTeXinterchartoks\germain@greek\@tempcnta=\expandafter{%
        \expandafter\germain@greekout\the\XeTeXinterchartoks\z@\@tempcnta}%
      \XeTeXinterchartoks\@tempcnta\germain@greek=\expandafter{%
        \the\XeTeXinterchartoks\@tempcnta\z@\germain@greekin}%
    \fi
    \ifnum\@tempcnta<4095\advance\@tempcnta\@ne\repeat}
% babel-french sets its own classes, and saves and restores the interchar
% state, each time a language is selected: join again after it does.
\addto\extrasfrench{\germain@greekjoin}
\addto\extrasenglish{\germain@greekjoin}
\AtBeginDocument{\germain@greekjoin}
\makeatother

\usepackage{geometry}
\geometry{a4paper,margin=25mm,marginparwidth=28mm}
\usepackage{marginnote}
\usepackage{xcolor}
\usepackage[normalem]{ulem}
\setlength{\emergencystretch}{3em}
\usepackage{hyperref}
\hypersetup{colorlinks=true,linkcolor=black,urlcolor=black,pdfborder={0 0 0}}

% --- Batch metadata ---------------------------------------------------------
% Taken as they stand by the rendering script, which shows them at the head of
% the reading view. They fix what the file claims to cover. The macro names are
% the parent project's, so that the scripts stay shared; \folder here means the
% volume's slug — `fr-9115` — and \pages counts Gallica views.

\newcommand{\folder}[1]{\gdef\grothFolder{#1}}
\newcommand{\batch}[1]{\gdef\grothBatch{#1}}
\newcommand{\pages}[2]{\gdef\grothFirst{#1}\gdef\grothLast{#2}}
\newcommand{\foldertitle}[1]{\gdef\grothFoldertitle{#1}}

% The holder's own shelfmark, printed as they write it: « Français 9115 ».
\newcommand{\shelfmark}[1]{\gdef\grothShelfmark{#1}}

% The dating, copied verbatim from the holder's catalogue — never inferred
% here. « XIXe siècle » is the BnF's; a pass that dates a leaf from its content
% says so in a \note{}, as its own inference.
\newcommand{\dating}[1]{\gdef\grothDating{#1}}

% The status notice, printed in the title block and repeated as a diagonal
% watermark on every page of the PDF. The manuscripts are in the public
% domain; what this line declares is not a right but a quality — an
% unverified first pass by a machine. The skills write the line; a file
% without it gets no watermark, so leaving it out is visible in the output.
\newcommand{\watermark}[1]{\gdef\grothWatermark{#1}}

\gdef\grothFolder{?}\gdef\grothBatch{}\gdef\grothFirst{?}\gdef\grothLast{?}%
\gdef\grothFoldertitle{}\gdef\grothShelfmark{}\gdef\grothDating{}\gdef\grothWatermark{}

% --- The résumé -------------------------------------------------------------
\newenvironment{resume}
  {\small\setlength{\parskip}{0.45em}}
  {\par\medskip\hrule\medskip}

% --- Keywords ---------------------------------------------------------------
% In English, closing the modernised reading's résumé: the modern vocabulary
% under which to search for what the leaves construct. The manifest extracts
% them to tag the volume — this line is the single source of the tags.
\newcommand{\keywords}[1]{\par\smallskip\noindent
  {\footnotesize\textsc{Keywords} --- \textit{#1}}\par}

% --- The critical apparatus -------------------------------------------------

% The Gallica view number, in the margin. This is the anchor that lets a
% disagreement over a formula be settled by looking at *one* image rather than
% a volume of 750. The site uses it to turn the facsimile as the transcription
% is scrolled. A view is one scanned image: usually one side of a leaf.
\newcommand{\page}[1]{%
  \marginnote{\footnotesize\color{gray!70}\textsf{v.\,#1}}%
  \label{page:#1}\ignorespaces}

% The BnF's pencilled foliation, where the leaf carries it and a pass can read
% it: « 198r ». This is what the literature cites — « ff. 198r–208v » — and
% the one line that turns a citation into a view. Recorded, never inferred:
% a folio a pass cannot read is not written.
\newcommand{\folio}[1]{%
  \marginnote{\footnotesize\color{gray!50}\textsf{f.\,#1}}\ignorespaces}

% The modernised reading's coarser counterpart to \page{}: which views a
% section is drawn from.
\newcommand{\pagerange}[2]{%
  \marginnote{\footnotesize\color{gray!70}\textsf{v.\,#1--#2}}%
  \ignorespaces}

% Illegible: never guessed.
\newcommand{\ill}{\textcolor{red!60!black}{[\,\ldots\,]}}

% A doubtful reading: the word is given, and flagged as uncertain.
\newcommand{\uncertain}[1]{\uline{#1}}

% An editorial addition — a missing letter, an implied word.
\newcommand{\add}[1]{\textcolor{blue!50!black}{[#1]}}

% Struck out by the writer. What was crossed out is part of the document.
\newcommand{\struck}[1]{\sout{#1}}

% The transcriber's note, on the state of the page or the mathematics.
\newcommand{\note}[1]{\par\smallskip{\small\color{gray!60!black}\textit{[#1]}}\par\smallskip}

% A marginal note of hers, distinct from the body of the text.
\newcommand{\marginal}[1]{\par\smallskip\noindent\hspace{1em}%
  {\small\color{gray!60!black}#1}\par\smallskip}

% --- Title ------------------------------------------------------------------
% The title block is French, like the documents themselves.

\AddToHook{shipout/background}{%
  \ifx\grothWatermark\empty\else
    \begin{tikzpicture}[remember picture, overlay]
      \node[rotate=52, scale=3.4, align=center, text=gray!18,
            font=\sffamily\bfseries]
        at (current page.center) {\grothWatermark};
    \end{tikzpicture}%
  \fi}

\AtBeginDocument{%
  \begin{center}
    {\large\bfseries Manuscrits de mathématiciens (Gallica) ---
      \ifx\grothShelfmark\empty\grothFolder\else\grothShelfmark\fi}\\[2pt]
    {\itshape\grothFoldertitle}\\[4pt]
    {\small vues \grothFirst--\grothLast
      \ifx\grothBatch\empty\else\ (lot \grothBatch)\fi}\\[2pt]
    \ifx\grothDating\empty\else
      {\small\color{gray!60!black}Datation du catalogue : \grothDating}\\[2pt]
    \fi
    \ifx\grothWatermark\empty\else
      {\small\bfseries\color{red!45!gray}\grothWatermark}\\[2pt]
    \fi
    {\footnotesize\color{gray!60!black}Transcription établie d'après les images IIIF de Gallica.
     Source gallica.bnf.fr / Bibliothèque nationale de France.\\
     Les lectures incertaines sont soulignées, les illisibles notées [\,\ldots\,],
     les ajouts entre crochets ; \textsf{v.} numérote les vues de Gallica,
     \textsf{f.} la foliotation au crayon de la BnF.}
  \end{center}
  \vspace{1em}}

\endinput


\folder{fr-9117}
\shelfmark{Français 9117}
\batch{5}
\pages{81}{100}
\dating{1801-1900}
\watermark{Lecture automatique\\non vérifiée}
\foldertitle{« Considérations générales sur l'état des sciences et des lettres aux différentes époques de leur culture, » par Sophie GERMAIN.}

\begin{document}

\note{Numérotation des feuillets. Chaque page porte en haut à gauche un
nombre à l'encre, biffé, de 71 (vue 81) à 90 (vue 100), qui semble une
ancienne pagination du manuscrit. Les rectos portent en haut à droite, au
crayon, la foliotation 40 à 49 (vues 81, 83, \ldots, 99). Les versos ne
portent pas de nombre au crayon, et aucun folio n'est porté pour eux.}

\page{81}\folio{40r}
les études étoient placés hors d'elle.

Nous avons cherché à imiter la littérature des anciens ; nous avons
adopté des fictions poëtiques qui ne se rattachoient plus, pour nous,
à des systêmes accrédités ou à des croyances adoptées. ces fictions,
jadis si riantes, étoient décolorées lorsqu'elles passoient dans les
écrits d'une nation qui ne les avoit pas imaginées ; leur signification
étoit énigmatique et conventionnelle pour nous, tandis \struck{étoient}
qu'elle\struck{s} étoi\struck{en}t pleine de sens et de vie entre les
mains des inventeurs. A une époque où l'imagination dominoit toutes
les conceptions humaines, les emblêmes dont nous nous efforçons de
faire revivre la grâce prêtoient un véritable secours aux conceptions
du poëte ; \struck{elles ornoient les} et un charme réel
\struck{prêtoient} à \struck{tous} \uncertain{ses} écrits \struck{un
charme réel}. Aujourd'hui les formes de style \struck{ff} sont plus en
harmonie avec notre caractère national. aussi voyons nous une école
nouvelle faire mille efforts pour créer une littérature qui nous soit
propre.

L'époque où nous vivons est remarquable par l'invasion des formes
mathématiques dans les ouvrages qui, par leur \struck{sujet} nature,
sont loin de pouvoir \struck{de présenter} atteindre l'exactitude
\struck{réelle} à laquelle de telles formes sont spécialement propres.
Il résulte de l'emploi maladroit des termes qui expriment une entière
certitude une sorte de déception intellectuelle dont la raison et le
goût sont également choqués. \struck{Une recherche de style, qui dont le
fond n'existe pas dans les ouvrages où l'on exprime}

\page{82}
Les personnes qui ne connoissent des sciences exactes que leurs
premiers élémens ont cru pouvoir reprocher aux géomètres une sécheresse
de style, qu'on a regardée comme inhérente au genre de leurs études.
Il est certain, cependant, que les ouvrages consacrés à l'exposition des
hautes théories mathématiques présentent dans le style même un charme
inexprimable. \struck{Il y a} On y \struck{trouve} remarque une
précision élégante, une extrême finesse, l'art de rendre présente à la
pensée une foule d'idées, qui pourtant ne sont pas textuellement
exprimées.

Tous ces avantages disparoissent d\struck{ans les} grotesques copies
qu'on introduit aujourd'hui dans le langage commun. on nous présente
\uncertain{hardiment} l'enveloppe sous laquelle nous sommes habitués à
trouver des pierres précieuses ; \struck{elle} et cette enveloppe
\struck{contient \ill{} a portant} des choses de peu de valeur. Nous
nous étonnons avec raison de \struck{les} voir dépourvues des ornemens
qui conviendroient au sujet. Pourquoi sont-elles astreintes aux
apparences de la solidité, tandis que celles de la légèreté seroient en
harmonie avec leur nature futile ? \struck{avec l'intérieur d'un tel
sujet ?}\note{Le « d » de « dans les » est laissé non biffé ; on attend
« des ». Sous « Nous », dans l'interligne, un petit mot lu avec doute
« Que ». « nature futile ? » est ajouté au-dessus de la ligne biffée,
dans laquelle « un » et « objets », surchargés, sont biffés aussi. Le
mot noirci au-dessus de « a portant » n'a pu être lu.}

Mais ce qu'il y a de plus vicieux c'est l'emploi des chiffres là où ils
n'expriment aucunes valeurs réelles ; ils usurpent le crédit dû aux
connoissances positives, et servent à établir l'erreur,

\page{83}\folio{41r}
\struck{à l'abri \ill{}} en donnant le change aux amis de la vérité.
\struck{Par Il} après avoir, renonc\struck{ant}é aux genres de littérature
dont les formes faciles et attrayantes sont propres à soutenir
l'attention, et qui d'ailleurs \struck{\ill{}} \uncertain{ont}
\struck{présentent} l'avantage d'être conformes à leurs habitudes, les
personnes qui, se \struck{sentant pénétrés} laissant entraîner par
l'amour des idées exactes, dont le besoin se fait actuellement sentir
plus généralement qu'à aucune autre époque, consentent à dévorer la
sécheresse attachée aux études élémentaires, mériteroient de trouver
dans les auteurs qui \struck{consacrent} leur servent de guide cette
conscience du vrai sans laquelle il est impossible d'atteindre aucuns
résultats importans.

Les nations éprouvent aujourd'hui le sort des individus qui se livrent
pour la première fois aux études sérieuses. Elles \struck{sont} encore
incapables de juger les ouvrages \uncertain{qui} \struck{elles veulent}
leur sont présentées ; elles s'indemnisent de la peine qu'elles prennent
à les étudier par une confiance aveugle dans les doctrines qu'elles y
trouvent et par le mépris des formes qu'avoient autrefois adoptées les
auteurs qui \struck{leur} promettoient une instruction moins solide.

Le pédantisme étoit jadis le défaut ordinaire des personnes adonnées
\struck{aux recherches} à l'étude ; maintenant, la plus obscure
médiocrité reléguée dans les provinces éloignées du centre

\page{84}
des mouvemens progressifs de la science offre seule quelques traces de
cet ancien défaut. mais, en revanche, \struck{ce} sont des masses
entières qui \struck{de} nous donnent le spectacle nouveau d'une
confiance illimitée dans leurs lumières et d'un mépris absolu pour les
personnes qui, fidèles à d'anciens documens, sont restées étrangères à
ce qu'on nomme aujourd'hui les nouvelles idées.

La jeunesse surtout renchérit sur cette ridicule manie, elle se croit
beaucoup trop instruite pour ne pas dédaigner le ton aimable de
plaisanterie qui chez notre nation accompagnoit autrefois une
instruction réelle et des opinions éclairées.

La littérature a besoin du secours des idées dominantes pour obtenir
l'attention de ses graves enfans. Il n'est pas permis de faire rire, si
ce n'est aux dépens des \struck{gens} personnes qui se montrent
ennemies des innovations. La raillerie est amère ; elle a perdu la
grâce qui savoit en adoucir les traits.

\note{Avant l'alinéa suivant, deux traits horizontaux à l'encre dans la
marge ; dans l'interligne, au crayon et d'une écriture qui paroît
autre, « Ch. 6 » suivi de mots à peine visibles, dont « Avenir de »
et \ill{} ; dans la marge de gauche, au crayon, « avenir de ».}

Ne nous alarmons pas de ces \struck{symptomes} symptômes : ils ne
seront que transitoires. Nous approchons de l'époque où le goût du
public pour les idées exactes déterminera le talent à s'occuper des
théories politiques. Lorsque la vérité aura trouvé des organes dignes
d'elle, elle paroîtra simple. Il sera alors facile de la reconnoître,
par sa nature, elle est étrangère à l'emphase, qui

\page{85}\folio{42r}
accompagne les doctrines sentencieuses de nos \struck{faux} pédans.
Bientôt on verra la politique recueillant le petit nombre de vérités qui
sont à son usage, adopter les formes qui conviennent à sa nature.
Agissant comme toutes les sciences qui ont besoin du secours de
l'expérience, elle craindra \struck{d'adopter} d'énoncer des théories
générales avant d'\struck{en} avoir justifié la réalité.

\marginal{Variante} [On verra alors ces progrès immenses dont on fait
tant de bruit se réduire à n'être autre chose que le développement
d'idées contenues dans les ouvrages de nos prédécesseurs. On les
trouvera à la vérité revêtues de formes \struck{plus appropriées}
nouvelles ; mais il sera clair que ces formes sont celles que les
sciences modernes ont adoptées, \struck{en sorte qu'après} Pendant un
tems \struck{où} une partie des connoissances humaines se distinguoit
des autres \struck{par} branches de la culture intellectuelle par une
méthode sévère et exacte, tandis qu'on retrouvoit partout ailleurs
\struck{le vague et \ill{}} des idées les plus heureuses unies
\struck{aux} conjectures les plus hazardées, traces des premiers essais
de la pensée ; après ce tems l'homogénéité qui fut le caractère des
travaux des anciens, dominés dans tous les genres par l'imagination
finira par se retrouver dans \struck{la culture} les travaux modernes
assujettis à la marche méthodique qui doit conduire à la connoissance
certaine des vérités propres à chaque genre d'étude.\note{Le mot « vérité »
de « à la vérité » est écrit dans la marge de gauche, en tête de ligne.
« unies », « après ce tems » et « dans tous les genres » sont ajoutés
dans l'interligne. Le crochet ouvrant est de l'auteur, en regard de la
mention marginale « Variante ».}

\page{86}
[Les lettres ont perdu de leur éclat ; elles n'attirent plus vers elles
l'attention des peuples ; elles ne sont plus l'objet de l'enthousiasme
de la jeunesse. La poësie, si elle ne trouve moyen de se rattacher à
quelques unes des idées qui intéressent les discussions politiques est
généralement dédaignée. Comment dans cette disposition des esprits
l'homme de génie pourroit-il trouver d'heureuses inspirations ?

Mais un jour plus pur ne tardera pas à renaître. L'analogie exige que
toutes les branches du savoir humain reçoivent des développemens, pour
ainsi dire, parallèles. L'attention se portera successivement sur
chacune d'elles, jusqu'à l'époque où, leurs progrès \struck{\ill{}}
devenant comparables, \struck{ou si j'ose le dire, identiques,} elles
obtiendront toutes ensemble le degré d'intérêt dû à leurs valeurs
respectives.

Ainsi voyons nous l'élève occupé d'acquérir les connoissances diverses
qui doivent composer l'éducation achevée, diriger son attention tantôt
vers \struck{les sciences tantôt vers l'étude des} un genre d'études,
tantôt \struck{vers} vers un autre genre tout différent du premier.
Chaque objet nouveau attire toutes les forces de son esprit, il semble
être devenu étranger à ce qu'il sait déjà, et on le croiroit incapable
d'aborder des questions qui \struck{pour} lui sont encore inconnues,
\struck{pour lui}. Cependant il arrive une époque, où chaque chose
reprend à \struck{aux} ses yeux \struck{de l'élève} son importance
véritable. Il voit des liaisons et des ressemblances là où il n'avoit
d'abord aperçu que des divisions et des différences. L'esprit humain
touche à une période semblable. Bientôt

\page{87}\folio{43r}
le tableau des sciences, des lettres, des beaux arts présentera au
spectateur \struck{à l'œil} attentif une symétrie méthodique qui lui
permettra d'embrasser d'un seul coup d'œil l'œuvre de l'esprit humain.
L'analogie qui a produit autrefois, pour les sciences des systêmes
hazardés et pour les lettres des allégories ingénieuses ou des
comparaisons pleines de grâces, prendra une force nouvelle. Elle ne
s'arrêtera plus à la superficie des choses, pour y chercher les
ressemblances qui s'offrent d'\struck{elles mêmes} à la première vue ;
elle pénètrera leur nature, et le type du vrai \struck{trouvera} offrira
dans les sujets les plus divers \struck{des copies variées de le guide le
plus} le caractère général de toutes connoissances certaines.

\struck{Laissons les sciences à part} Nous avons traité plus haut de la
révolution qui s'est opérée dans la manière dont nous envisageons les
sciences \struck{math} physiques ; nous avons dit comment les méthodes
géométriques ont étendu leur empire en portant la certitude qui leur est
propre dans des régions qui furent longtems le domaine des idées
systématiques. Les sciences morales et politiques ne tarderont pas à
subir la même transformation ; déjà l'opinion publique s'attend à ce
changement, et en devance même la réalisation par un enthousiasme
irréfléchi pour les doctrines qui lui en offrent l'espérance. Cette
erreur passagère est exempte de danger. elle sera peu durable et le goût
dont elle est le symptôme sera bientôt pleinement satisfait.\note{Dans la
marge de droite, en regard des dernières lignes, de courts traits
verticaux à l'encre, peut-être le crochet fermant des passages ouverts
aux vues 85 et 86 ; rien ne permet de l'assurer.}

\page{88}
les méthodes existent ; une difficulté née de l'amour propre peut seule
en retarder l'emploi. Les hommes en état de traiter de pareilles
questions \struck{ne jugent pas qu} craignent de n'être pas estimés de
leurs pairs, et de ne pas rencontrer de juges éclairés dans les
personnes étrangères aux sciences.

Un pareil obstacle ne peut pas subsister \struck{bien} longtems : et
nous pouvons \struck{donc} regarder dès à présent les sciences morales
et politiques comme appartenant au domaine des idées exactes.

Mais l'analogie se montrera dans tout son charme et dans toute sa
puissance \struck{se montrera} lorsque l'esprit d'examen entreprendra de
comparer la manière d'agir des lettres et des arts, \struck{et} lors que
portant ensuite ses regards vers les modèles offerts par le spectacle de
la nature, il trouvera de tout côté des copies sans cesse renouvelées de
ce modèle du vrai qui appartient à notre être ; qui est la source de
toutes nos plaisirs intellectuels et qui, réfléchi autour de nous,
\struck{produit à la fois} devient la cause des impressions que nous
recevons des objets \struck{qui} dont est frappée notre imagination,
\struck{et le principe de l'imitation, en vertu à l'aide sert à la fois
duquel} le génie \struck{sait nous émouvoir en} produi\struck{sant}t à
son gré des émotions \struck{dont} qui se ressemblent et sont dues
cependant à des causes qui diffèrent comme la nature des moyens
\struck{qu'il sait employer \ill{} que constituans} constitutifs
d'autant d'arts séparés, \struck{donne immédiatement du sentiment} Il
doit cette puissance au principe de l'imitation et l'imitation dérive
du sentiment d'analogie. S'il n'existoit pour nous aucun type commun
entre les divers objets dont nous recevons l'impression, les arts
auroient pu copier \struck{ce qui}\note{Page très corrigée. « née de »,
« et » (devant « des arts »), « lors que », « ensuite », « devient la
cause des », « dont est », « constitutifs », « sont » sont écrits dans
l'interligne ; « en vertu » et « à l'aide », eux-mêmes biffés, aussi.
« se » est ajouté dans la marge de gauche devant « ressemblent », écrit
en surcharge d'un premier mot. Les deux lignes « Il doit cette
puissance au principe de » et « l'imitation et l'imitation dérive du
sentiment » sont écrites au-dessus et au-dessous de la ligne biffée
« donne immédiatement du sentiment ». L'ordre de lecture retenu ici est
celui qu'impose le sens.}

\page{89}\folio{44r}
les objets extérieurs, les lettres auroient pu redire les événemens
dignes d'attention ; mais ni les uns ni les autres n'auroient su
reproduire une impression semblable à celle qui vient des choses
existantes en employant \struck{des choses existantes} des moyens qui ne
leur auroient pas été immédiatement fournis par le sujet.\note{Un long
trait à l'encre, en boucle, entoure le passage de « reproduire » à
« sujet » ; « existantes » est écrit dans l'interligne. En tête de la
page, au crayon et d'une écriture qui paroît autre, un mot lu avec
doute « Double ».}

\struck{C'est à l'observation} Les lois de l'être établissent de
certains rapports entre un \uncertain{module} donné et tout ce qui tient
au sujet auquel ce \uncertain{module} appartient. \uncertain{Ce sont}
ces rapports qui agissent sur notre imagination. L'âme peut être affectée
de la même manière par l'entremise de sens différens, parce qu'elle
reçoit alors le même ordre d'idées. ainsi l'éloquence, l'art musical et
la peinture \struck{procèdent d'une manière} \uncertain{pantomime}
procèdent d'une manière analogue ; la peinture ne dispose que d'un seul
instant ; mais elle sait le choisir de façon à rappeler ceux qui ont
précédé, et à faire pressentir la suite de l'action que le tableau
représente.

Si on excepte un petit nombre de cas dans lesquels l'auditeur est
préparé, par les circonstances du moment, à prêter toute son attention à
l'orateur, celui-ci commence son discours avec calme ; le musicien agit
de même au commencement de sa pièce ; il emploie une modulation simple,
et \struck{le même autre} réserve \struck{le même} aussi les mouvemens
expressifs pour la suite de l'action. Dans le \struck{nat} spectacle de
la nature ;\note{« aussi » est écrit dans la marge de droite. La page
s'achève sur un long trait horizontal après « nature ; ».}

\page{90}
nous retrouvons la même gradation. le lever du soleil est précédé d'un
crépuscule : et un autre crépuscule \struck{termine} annonce la fin de
son apparition. \struck{Le goût exige que} aussi voyons nous l'écrivain
et l'orateur termine\struck{nt}r leur oraison avec la même simplicité
qu'elle \struck{\ill{}} a été commencée. \uncertain{Une pensée} d'éclat
placée à la fin d'un discours laisseroit l'auditeur dans une sorte
d'éblouissement qui seroit fatiguant\struck{. \ill{}} : \struck{cette
manière est repoussée} le goût, en s'épurant a proscrit cette manière.
L'influence de l'analogie est tellement sympathique que le musicien
renonce aussi à nous faire entendre les bruyans accords par lesquels il
étoit naguère d'usage de terminer l'acte de cadence. Dans les morceaux
\struck{plus} modernes, \struck{la tonique arrive sans éclat et
l'oreille est ramenée au repos absolu par le des sons par} les dernières
mesures préparent le repos absolu ; non pas seulement par la modulation
qui doit amener la tonique, mais encore par \struck{la} diminution
d'intensité des sons \struck{a} employés vers la fin du morceau
\struck{dans les dernières mesures}.\note{« annonce la fin » et « aussi
voyons nous » sont écrits dans l'interligne. Au-dessus de « qu'elle »,
quelques mots biffés et noircis. Devant « d'éclat », un mot noyé sous
une tache d'encre ; au-dessous, un mot biffé. « vers la fin du
morceau », écrit au-dessus de « dans les dernières mesures », porte
lui-même un trait qui pourroit être une rature.}

[Nous allons voir combien il existe de ressemblance véritable entre les
manières \struck{de faire} pratiquées \struck{employées} \struck{par les}
dans les arts, qu'on a coutume de regarder comme \struck{étant}
entièrement étrangers l'un à l'autre.

[Ainsi en comparant avec plus de détails \struck{l'art dont} les grands
effets produits \struck{par} d'un côté par le talent de l'orateur, de
l'autre par celui de l'habile compositeur, nous aurons occasion de nous
convaincre \struck{à quel point}

\page{91}\folio{45r}
de l'identité des rapports qui agissent sur notre imagination. Une seule
condition est nécessaire pour les appercevoir : il faut être familier
avec chacun des \uncertain{modules} qui seroient de nature commune à de
tels rapports.

\struck{Dans} Les effets d'une grande puissance frappent à la vérité les
hommes les moins instruits ; mais la finesse du tact a besoin
d'exercice. le goût, l'oreille sont susceptibles de se former ; et
souvent on se croit absolument incapable d'apprécier \struck{les effets}
l'harmonie \struck{musicale} à cause du préjugé qui sépare la musique du
domaine de l'intelligence. Au contraire, personne ne veut paroître
insensible aux beautés littéraires : et une éducation à laquelle on
attache une \struck{si} haute importance, prépare l'homme du monde à
porter des jugemens raisonnables, ou au moins à savoir choisir les
autorités qui \struck{doivent} \struck{suppléer à} lui fournir les
opinions qu'il pourra émettre sans honte. A l'égard de la musique les
choses se passent autrement ; l'enfant qui ne \struck{démontre} trouve
pas une grande sensibilité aux \struck{par} premiers accords qu'on lui
fait entendre est regardé comme n'étant pas organisé pour cet art ;
souvent l'enseignement est \struck{souvent} mauvais ; \struck{et} l'on
juge, d'après le défaut de progrès, que l'élève est incapable
d'apprendre, et \struck{et} l'argument tiré de l'inutilité de l'art en
lui même, fait \struck{\ill{}} qu'on l'abandonne bientôt. Ainsi, à moins
d'être né assez heureusement pour annoncer, dès les premiers essais,
des\note{« vérité » (dans « à la vérité ») est écrit dans la marge de
gauche ; « Au contraire », « ne », « trouve », « souvent » et « et »
sont ajoutés dans l'interligne.}

\page{92}
dispositions remarquables, si la famille est elle même étrangère au goût
de la musique, \struck{l'enfant} l'enfant se trouve privé des secours qui
l'auroient pu conduire à la connoissance des beautés en ce genre. Il est
évident que, si on agissoit de la même manière par rapport aux lettres,
les hommes doués d'une manière de sentir \struck{tellement vive qu'elle
n'attend pas le jugement des autres hommes pour} exquise seroient les
seuls qui \struck{fussent} acquerreroient la connoissance des chefs
d'œuvres littéraires. Je crois qu'il est également rare de voir l'enfant
annoncer de grandes dispositions dans l'un ou l'autre genre ;
l'éducation et l'opinion rendent compte de la différence entre les
nombres qui expriment combien il y a d'homme capables de bien juger en
littérature et combien il s'en trouve qui sauront apprécier le mérite
d'une composition musicale. \uncertain{Chévenon}\note{Mot d'une écriture
plus grande et plus appuyée, en fin de ligne, souligné d'un trait
distinct : peut-être un nom propre, ou un renvoi ; la lecture est
douteuse. « l'enfant », répété dans l'interligne au-dessus du mot biffé,
et « exquise » sont des additions.}

On croit avoir trouvé une objection fondamentale dans la disposition
naturelle nécessaire pour juger la justesse d'un son. Mais, lorsqu'il ne
seroit pas certain que l'oreille apprend à distinguer les tons entre
lesquels l'oreille ne trouvoit d'abord aucune différence, il resteroit
encore une foule de choses, qui, étant entièrement indépendantes du
choix du \uncertain{module}, pourroient être appréciés par l'homme doué
d'intelligence et accoutumé à porter son attention sur les effets de la
musique.

\page{93}\folio{46r}
\struck{On s'étonne aujourd'hui de ce que les anciens ont attribué au choix
des genres une importance dont les raisons échapent \struck{aux} à notre
attention. Cette haute importance l'influence qu'avoit eu le choix des
modes semblent attester que les grecs ont en effet regardé la musique
comme une véritable langue et lorsque nous aurons}\note{Paragraphe barré
d'une grande croix à l'encre, et resté inachevé.}

On ne comprend plus aujourd'hui ce que l'histoire nous apprend touchant
l'influence des différens modes ; et comment parviendroit-on à s'en
rendre compte, lorsqu'on s'obstine à regarder la musique comme l'art de
flatter l'oreille ? réduite à cet unique usage, pourroit elle
\struck{encore l'ins} être l'objet d'une attention sérieuse ? auroit elle
pu produire les effets \struck{mer} que les anciens nous racontent ?

Mais ces merveilles cesseront de nous étonner lorsqu'en comparant les
moyens usités en musique à ceux que l'orateur sait employer, nous aurons
mis dans tout son jour cette vérité, reconnue maintenant par un petit
nombre de personnes, qui craignent même de l'énoncer, que la musique est
une langue, et une langue énergique. Elle emploie les sons ; mais les
sons ne la constituent pas. Elle a ses phrases, ses périodes, ses repos,
ses hardiesses. Ses beautés flattent l'oreille ; \struck{mais} mais
\struck{et} ne s'y arrêtent pas ; elles pénètrent l'âme et peuvent
exercer sur elle un empire véritable. Ainsi la poësie emploie des
\struck{mots qui} sons articulés agréables à l'oreille, et l'on auroit
une idée fort incomplète du charme qui y est attaché si on oublioit le
sens des phrases, pour ne \struck{s'attacher} s'occuper que
\struck{de} leur nombre.

\page{94}
La musique est toute métaphysique ; ses expressions sont générales ; elle
ne possède aucun nominatif ; elle ne peut exprimer que des sentimens ;
mais il est en sa puissance de mettre l'âme de l'auditeur dans la même
situation que \struck{si elle} l'a occupée d'un récit positif d'une
action particulière.\note{« l'a », écrit au-dessus de « si elle » biffé,
est lu avec doute.}

Cette langue procède comme toutes les autres langues ; elle
\struck{aussi} fournit des expressions de tous genres.

Le compositeur doit, comme l'orateur, avoir une idée dominante. Il
s'empare d'abord de l'attention de l'auditeur en ne lui présentant que
des phrases usitées, \struck{des idées dont tout le monde tombe d'accord}
que personne ne puisse s'étonner d'entendre. Elles amèneront le
développement du sujet ; mais à moins de circonstances particulières,
elles ne \struck{le contiendront} l'expliqueront pas en entier.

Si le littérateur veut occuper son lecteur d'idées gracieuses et
\struck{bien} soutenir l'attention sans produire aucune impression qui
puisse troubler le repos de l'âme, ses phrases \struck{seront} seront de
l'harmonie ; il \struck{évitera} rejètera toute expression ambitieuse ;
il sera constamment pur, et jamais recherché.

Le compositeur qui se propose le même but employera des moyens tout
pareils ; sa composition sera toujours correcte ; elle sera simple ;
l'oreille ne s'étonnera d'aucune des phrases qui lui seront offertes ;
mais un charme \struck{toujours}

\page{95}\folio{47r}
\struck{sont toujours} soutenu fera pénétrer dans l'âme le sentiment de la
grâce.\note{Suite de la dernière ligne de la vue 94 : « mais un charme
\ldots{} soutenu ».}

L'homme de lettre veut-il composer un ouvrage propre à égayer le lecteur,
il connoîtra l'art d'amener les contrastes ; mais s'il craint de tomber
dans le burlesque, il évitera les tours disparates et les expressions
\struck{propres} qui conviendroient à la peinture d'un sujet tragique.
L'orateur se propose \struck{des objets} un but plus important que
d'entretenir l'âme de ses auditeurs dans une situation douce et calme ou
de faire \struck{passer} naître la gaîté autour de lui ; mais la
conversation des \struck{gens} personnes aimables n'a pas \struck{d'autres}
besoin d'autres \struck{secours} effets pour nous plaire et nous attacher.

Elles peuvent nous donner lieu de remarquer que la plus part des choses
qui \struck{ont} ont le caractère de la grâce \struck{si elles sont}
dites d'une certaine manière, \struck{peuvent} \uncertain{deviendroient}
\struck{propres} au moyen \struck{à l'aide} de foibles changemens et
d'une autre manière de les dire, propres à prendre le caractère de la
gaîté.\note{Au-dessus de « peuvent », dans « Elles peuvent », le mot
« conversations » ; dans la marge, « de » souligné et une croix qui
renvoie à « Elles » : l'auteur semble avoir voulu « Ces conversations
peuvent ». Un trait à l'encre entoure « ont le caractère de la grâce » et
« d'une certaine manière ». « deviendroient » est écrit en surcharge
au-dessus de mots biffés et n'est lu qu'avec doute ; « propres » et « de »
(devant « la gaîté ») sont ajoutés dans l'interligne.}

[Le rapport entre la grâce et la gaîté \struck{devient} est sensible dans
les compositions musicales.

Les mêmes modes, les mêmes coupures de phrases, y \struck{conviennent}
sont employées dans l'un et l'autre genres ; \struck{la} et le changement
de mouvement tient la place de celui des manières de dire.

\page{96}
Plus de vitesse suffit pour donner le caractère de la gaîté à une
composition, qui, \struck{dite avec} exécutée plus lentement eut
\struck{\ill{}} été simplement gracieuse. [Mais, en musique comme en
littérature, \struck{celle-ci exclut} la grâce, exclut les contrastes,
qui conviennent à la gaîté ; comme la gaîté, sous peine de tomber dans
le burlesque, rejette le \struck{aussi les contrastes, qui eussent été
bien placés dans l'autre ; le burlesque résulteroit de l'emploi d'un
autre} genre de contrastes réservé à l'expression des grandes
émotions.\note{« exécutée » et les deux lignes « la grâce, exclut
\ldots{} sous peine » et « de tomber dans le burlesque, rejette le » sont
écrits dans l'interligne ; au-dessus de « en musique comme en
littérature », une ligne entièrement noircie n'a pu être lue. L'ordre
retenu est celui qu'impose le sens.}

\struck{Nous venons de parler des} Les rapports que nous venons de
remarquer entre deux situations de l'âme où elle est dans un état de
bien être \struck{se font remarquer} \uncertain{également sentis} entre
celles où l'âme éprouve un sentiment de tristesse.\note{« également
sentis » est écrit au-dessus de « font remarquer » biffé ; « sentis »
paroît biffé à son tour.}

La mélancolie et le désespoir peuvent être produits par une seule et
même cause. Il arrive même que ces deux manières d'être affectée se
succèdent alternement, \struck{et se reproduisent} chez le même sujet
jusqu'à ce que, \struck{peu} à peu l'impression reçue s'affoiblissant
peut éloigner ou même fasse disparoître l'accent du désespoir, pour ne
plus conserver que l'impression \struck{plus douce} d'une tristesse
\struck{qui} \uncertain{qu'un} susceptible de distraction. \struck{et}
Alors la mélancolie n'a plus le caractère de l'abattement ; elle devient
douce, et quelque fois chère aux personnes qui en sont
affectées.\note{La phrase « jusqu'à ce que \ldots{} de distraction »
est laissée telle qu'elle est écrite, avec « peut éloigner ou même
fasse » et « qu'un susceptible », qui ne se construisent pas. Une plume
posée sur la page a laissé une trace en travers de « et se
reproduisent », mot qui est aussi biffé.}

\page{97}\folio{48r}
les différens genres, liés entre eux par une origine commune, inspirent à
l'homme de lettres, au poëte et à l'orateur des compositions qui
\struck{à l'imitation} ont aussi entre elles des rapports sensibles, et
pourtant des caractères divers.\note{La vue 96 s'achève sur une phrase
complète (« qui en sont affectées. ») et celle-ci s'ouvre sur « les
différens genres » sans majuscule : le raccord entre les deux pages ne
se fait pas.}

La sombre mélancolie exprime les mêmes idées que le désespoir ; dans la
simple conversation, \struck{il suffit de changer le ton mouvement
ralentir ou précipiter le discours pour} le passage de l'un à l'autre peut
être uniquement marqué par le ralentissement \struck{et} par la
précipitation du discours. — La même chose est vrai par rapport aux
compositions musicales de peu d'étendue. \struck{une} complainte
\struck{où} \struck{insp} dans laquelle le caractère d'une sombre
tristesse se trouveroit fortement exprimé, feroit entendre les accens du
désespoir si on en précipitoit le mouvement.\note{« complainte » est
laissé sans article, l'article « une » ayant été biffé.}

[En littérature, aussi bien que dans l'art musical, les choses se passent
autrement lorsqu'il s'agit de compositions plus importantes. Alors les
caractères sont distincts ; et une bonne composition dans un des genres
\struck{pourroit} deviendroit ou foible ou tout à fait mauvaise si on se
contentoit de changer la manière de les dire.\note{Un trait à l'encre
entoure « Alors les caractères sont distincts ».}

\page{98}
[Il est facile de comprendre la raison de cette différence de facture
entre les grandes compositions et celles qui ont peu d'étendue, et, je
pourrois dire peu d'importance sous le rapport de l'art.

Les sentimens douloureux qui produisent les deux affections dont nous
examinons l'expression \struck{admettent pendant} se manifestent
également par l'une et par l'autre ; \struck{il est dans la nature} notre
manière de sentir ne nous permet pas de demeurer longtems dans la
situation violente qui caractérise le désespoir. Cette situation
déchirante amène une fatigue extrême, et, si elle n'est \struck{pas} pas
portée à un \struck{tel} degré d'exaltation assez fort pour troubler
entièrement nos facultés intellectuelles, la fatigue \struck{ne} forcera
notre âme à prendre quelque repos. Ce genre de repos n'est pas celui du
bien être ; il approche de la stupidité ; c'est l'abattement de la
douleur ; c'est une mélancolie sombre, \struck{fait place alors c'est}
une tristesse profonde.\note{Le second « c'est » est écrit dans
l'interligne, au-dessus de « une mélancolie ».}

\page{99}\folio{49r}
\note{En tête de la page, au-dessus de la première ligne, le mot
« Georges », à l'encre, entouré d'un paraphe, d'une écriture qui paroît
autre que celle du texte.}

[Quel que soit le genre de composition destiné à reproduire de telles
impressions, il est assujetti à se conformer aux besoins de notre
sensibilité ; il doit donc employer \struck{concurremment}
alternativement la peinture de l'extrême \struck{tristesse} tristesse
\struck{il résulte de la} et celle du désespoir. ainsi ces deux aspects
différens d'un seul et même sentiment se trouvent pour ainsi dire
opposés l'un à l'autre, et dans un tel degré de rapprochement que la
comparaison en devient inévitable. Il convient alors d'établir entre
leur expression \struck{toutes} \struck{les} d'autres différences
\struck{possibles} que celles du mouvement. La nature de ces affections
en offre le moyen.\note{« donc », « alternativement », « d'autres » et
« que celles du mouvement » sont écrits dans l'interligne.}

Non seulement l'abattement s'exprime avec lenteur ; mais une sorte de
monotonie lui convient ; le poëte, l'orateur, choisiront des syllabes
douces, faciles à prononcer, et dont l'émission semble \struck{devoir}
exiger peu d'efforts. Le musicien aura recours à des notes
\struck{de presque} de même valeur ; en sorte que les sons qu'il fait
entendre produisent sur l'âme \struck{un} un effet \struck{le plus}
voisin du silence absolu, qui \struck{produiront une} \uncertain{mieux}
\uncertain{nourrit} \uncertain{mieux} la préoccupation dont il doit nous
faire connoître l'existence.\note{Fin de phrase très remaniée : au-dessus
de « qui produiront une préoccupation », biffés en partie, sont écrits
trois mots, lus avec doute « mieux », « nourrit », « mieux » — le
troisième pourroit être « malheur » —, et « la » ; « sur l'âme » est
dans la marge de gauche, « un » et « même » dans l'interligne.}

\page{100}
[L'attention à suivre ces règles offre le double \struck{effet} avantage
de renforcer l'effet qui seroit produit par le seul changement de la
manière de dire et de préparer une opposition plus frappante entre les
accens de la profonde mélancolie et ceux du désespoir.

De même aussi les affections violentes, lorsqu'elles devront se montrer
dans une même composition à côté de la \struck{douleur profonde} et
sombre douleur, trouveront d'autres nuances que celles qui résulteroient
de la seule précipitation du mouvement. Des expressions fortes et
énergiques, des notes dures à l'oreille avertiront l'âme de l'exaltation
nouvelle produite par le sujet.\note{« et sombre douleur » : « et » est
écrit après les mots biffés, « douleur » en tête de la ligne suivante,
avec un renvoi.}

Le poëte, l'orateur, l'auteur dramatique, littérateur ou musicien, tirent
\uncertain{ses} grands effets de l'art avec lequel ils savent amener un
mot une note inattendus. L'âme de l'auditeur s'étoit identifiée avec le
développement d'une action qui lui étoit, pour ainsi dire, présente ;
tout à coup, elle voit avec surprise un incident nouveau qui en renforce
\struck{trop} l'impression ; elle se trouble, \uncertain{alarmée} d'un
surcroît de malheur, dont elle ne peut plus mesurer
l'étendue.\note{« tirent » et « ses » sont écrits en surcharge ; la
finale de « inattendus » est surchargée elle aussi. « alarmée », en
petits caractères dans l'interligne, est lu avec doute.}

\end{document}
